% !TEX root = ../thesis.tex

\chapter{Background and Related Work}
\label{ch:background}

% This chapter migrates and expands proposal §4 (Literature Review) and adds
% the technical background sections needed to support Chapters 4 and 5.

\section{Medical Tabular Data and Privacy Regulations}
\label{sec:bg:regulations}

Classical anonymisation methods like generalisation, suppression, distortion,
swapping, masking, and the $k$-anonymity / $\ell$-diversity / $t$-closeness
lineage struggle to preserve utility on high-dimensional EHR data and have
repeatedly been shown vulnerable to linkage and inference
attacks~\cite{anonymization,Murtaza2023}. Synthetic data sidesteps direct
release of real records but does not automatically resolve the trade-off:
\textcite{stadler2022synthetic} show that non-DP synthetic data still leaks
information about outliers, motivating the formal privacy machinery
(§\ref{sec:bg:dp}) used by the DP generators in this thesis. 
\section{Synthetic Data Generation: A Taxonomy}
\label{sec:bg:taxonomy}

Synthetic tabular data generators fall broadly into statistical models (copulas,
Bayesian networks, sequential trees), deep generative models (GAN- and
VAE-based), and language-model-based approaches that treat rows as
text~\cite{hernandez2022synthetic,Murtaza2023}. Alongside these data-driven
families, \textcite{Murtaza2023} also describe a knowledge-driven approach that
builds generators from hand-curated expert rules, and a hybrid approach that
blends the two; both remain labour-intensive and predominantly knowledge-driven,
and therefore fall outside the fully data-driven scope of this thesis. Each
family admits a differentially-private variant, yielding the GAN/LLM $\times$ DP/non-DP axes
along which the five generators evaluated in this thesis are positioned. The
remainder of the chapter develops the two families used here
(§\ref{sec:bg:gan}, §\ref{sec:bg:llm}) and the DP machinery they rely on
(§\ref{sec:bg:dp}). Unlike image 
or text data, which consist of uniform basic units amenable to a single modelling 
paradigm, tabular data combines categorical and numerical columns with heterogeneous 
statistical properties which is a challenge that no single generative family solves cleanly, 
and that drives the proliferation of distinct approaches surveyed in \cite{shi2025comprehensive}.
 Within each family, known failure modes (training instability and mode collapse in GANs,
 a tendency toward over-averaged outputs in VAEs, and hallucination in LLMs) have
 motivated successive specialisations, including the diffusion-model and fine-tuned-LLM
 branches that occupy the DP/non-DP axes evaluated in this thesis.


\section{GAN-based Approaches for Tabular Data}
\label{sec:bg:gan}

Generative adversarial networks constitute the dominant deep generative
paradigm for tabular synthesis, yet inherit the training instability and
mode-collapse pathologies first documented in the image
domain~\cite{shi2025comprehensive}. The remainder of this section reviews
the architectural progression from the original minimax formulation through
the Wasserstein and gradient-penalty refinements that stabilise adversarial
training (§\ref{sec:bg:gan:foundations}), and subsequently examines the two
healthcare-oriented variants situated on the DP and non-DP axes evaluated
in this thesis: HealthGAN (§\ref{sec:bg:gan:health}) and
PATE-GAN (§\ref{sec:bg:gan:pategan}).


\subsection{Foundations: GAN, WGAN, WGAN-GP}
\label{sec:bg:gan:foundations}

Generative adversarial networks formulate data generation as a
minimax game between a generator $G$, which maps noise samples into
the data space, and a discriminator $D$, which estimates the
probability that a given sample originates from the data distribution
rather than from $G$~\cite{Goodfellow2014GAN}:
\begin{equation}
\label{eq:gan-minimax}
\min_{G} \max_{D}\; V(D, G) =
\mathbb{E}_{x \sim p_{\text{data}}}\!\left[\log D(x)\right]
+ \mathbb{E}_{z \sim p_z}\!\left[\log\bigl(1 - D(G(z))\bigr)\right].
\end{equation}
The global optimum is attained when the generator distribution
matches the data distribution and the discriminator is everywhere
indifferent. In practice, however, optimisation of this objective is
fragile: the discriminator can saturate and starve the generator of
gradient signal, the value function offers no monotonic indicator of
convergence, and the generator is prone to mode collapse, in which
it concentrates mass on a narrow subset of plausible samples.

These pathologies have been attributed to the implicit
Jensen--Shannon divergence underlying the original objective, which
is discontinuous in the generator's parameters whenever the model
and data distributions have disjoint or near-disjoint
supports~\cite{Arjovsky2017WGAN}. Replacing it with the Wasserstein-1
distance restores continuity throughout training and yields a
gradient that remains informative even in this regime. Via the
Kantorovich--Rubinstein duality, the Wasserstein-1 distance admits
the variational form
\begin{equation}
\label{eq:wasserstein-dual}
W(p_{\text{data}}, p_g)
= \sup_{\|f\|_{L} \leq 1}\;
\mathbb{E}_{x \sim p_{\text{data}}}\!\left[f(x)\right]
- \mathbb{E}_{x \sim p_g}\!\left[f(x)\right],
\end{equation}
so that the discriminator is replaced by a 1-Lipschitz \emph{critic}
whose scalar output approximates this supremum. The Lipschitz
constraint is enforced in the original Wasserstein GAN by clipping
the critic's weights to a fixed interval, a modification that
produces markedly stabler training and a loss that correlates with
sample quality.

Weight clipping, however, was subsequently shown to be the source of
its own pathologies, chief among them a sharp loss of critic
capacity and gradient norms that either explode or vanish depending
on the clipping threshold~\cite{Gulrajani2017WGANGP}. The improved
formulation replaces clipping with a soft penalty on the critic's
input-gradient norm,
\begin{equation}
\label{eq:wgan-gp}
\mathcal{L}
= \mathbb{E}_{\tilde{x} \sim p_g}\!\left[D(\tilde{x})\right]
- \mathbb{E}_{x \sim p_{\text{data}}}\!\left[D(x)\right]
+ \lambda \,
\mathbb{E}_{\hat{x}}\!\left[
\bigl(\|\nabla_{\hat{x}} D(\hat{x})\|_{2} - 1\bigr)^{2}
\right],
\end{equation}
where $\hat{x}$ is sampled uniformly along straight lines between
paired real and generated samples. The penalty encourages the
critic to satisfy the Lipschitz condition with unit norm along the
region of the input space where the two distributions are being
compared, rather than uniformly across its parameters. This
Wasserstein-with-gradient-penalty formulation is the backbone on
which HealthGAN (§4.2) and its differentially-private extension
(§2.5) are built.

\subsection{Healthcare-specific Variants: HealthGAN}
\label{sec:bg:gan:health}

HealthGAN was introduced as a response to the limitations of medGAN, which
could only generate binary diagnosis codes and often produced unrealistic
distributions, making it unsuitable for mixed-type EHR
data~\cite{healthgan,yale2020synthesizing,bhanot2022investigating}. To address
these shortcomings, the model combines medGAN's design with the WGAN-GP
framework, leveraging the Wasserstein distance and a tailored data
transformation pipeline to support continuous and categorical variables
commonly found in datasets such as MIMIC-III~\cite{healthgan}. Its architecture
employs large batch sizes to ensure that outliers and rare clinical values are
reliably captured, and incorporates a bottleneck to reduce memorisation and
maintain a compact model footprint, facilitating secure model export outside
the training environment. A key contribution of HealthGAN is the introduction
of \emph{Nearest-Neighbour Adversarial Accuracy} (AA), a metric suite that
enables joint assessment of resemblance and privacy and has been recognised in
systematic reviews as an important development in synthetic health data
evaluation~\cite{Murtaza2023}. Empirical results show that HealthGAN was the
only model in a six-method benchmark that maintained privacy while supporting
model export, achieving competitive utility and resemblance relative to
established baselines~\cite{healthgan}.

\subsection{PATE-GAN}
\label{sec:bg:gan:pategan}

\textsc{PATE-GAN} adapts the Private Aggregation of Teacher Ensembles (PATE)
framework to generative modelling by training multiple teacher discriminators
on disjoint data partitions and aggregating their predictions through a
differentially-private noisy mechanism, thereby providing formal
$(\varepsilon,\delta)$-DP guarantees while avoiding direct exposure of private
gradients~\cite{pategan,Papernot2017PATE}. This provides a tunable privacy
budget with an explicit accountant, which aligns with survey guidance to
report both formal guarantees and empirical privacy
audits~\cite{hernandez2022synthetic,Murtaza2023}. In domains where non-DP
synthetic data has been shown to leave outliers
vulnerable~\cite{stadler2022synthetic}, and where DP graphical approaches
(e.g., PrivSyn~\cite{privsyn2021}) trade utility for stronger privacy,
PATE-GAN offers a principled middle ground with controllable $\varepsilon$.

\section{LLM-based Approaches for Tabular Data}
\label{sec:bg:llm}

Large language models offer a generative paradigm that is structurally
distinct from the adversarial approach of \S\ref{sec:bg:gan}. Rather than
learning an implicit sampler through a minimax game, a decoder-only
transformer models the joint distribution of a record explicitly, as a
product of autoregressive conditionals over a tokenised representation of
that record. This formulation confers two properties that are attractive for
heterogeneous clinical tables. First, mixed-type columns are handled without
the mode-specific normalisation or one-hot expansion required by conditional
tabular GANs, because every value, whether numeric or categorical, is
rendered as text~\cite{Miletic2024}. Second, the self-attention mechanism
models dependencies across all columns natively, which mitigates the
difficulty GANs encounter as dimensionality grows~\cite{Miletic2024,shi2025comprehensive}.
The remainder of this section describes the row-as-text serialisation that
makes tabular data amenable to language modelling
(\S\ref{sec:bg:llm:tabula}), the Pythia model family used as the synthesiser
(\S\ref{sec:bg:llm:pythia}), the low-rank adaptation scheme used for
parameter-efficient fine-tuning (\S\ref{sec:bg:llm:lora}), and the
differentially private regime within which this thesis situates the
Pythia-1B-DP variant (\S\ref{sec:bg:llm:dp}).

\subsection{Tabula and the Row-as-Text Paradigm}
\label{sec:bg:llm:tabula}

The serialisation underlying language-model-based tabular synthesis was
introduced by the GReaT framework, which encodes each record as a
natural-language sentence~\cite{Borisov2023GReaT}. A row with features
$f_1, \dots, f_m$ and corresponding values $v_1, \dots, v_m$ is rendered as a
sequence of clauses of the form ``$f_i$ is $v_i$'', concatenated into a single
token sequence $s = (w_1, \dots, w_T)$. During training a random permutation
of the feature order is applied to each row, so that the model is not biased
towards a fixed column ordering and can subsequently be conditioned on any
subset of features at sampling time~\cite{Borisov2023GReaT}. The language
model is then fit with the standard autoregressive objective, which
factorises the sequence probability as
\begin{equation}
\label{eq:bg:llm:ar}
p_\theta(s) = \prod_{t=1}^{T} p_\theta\!\left(w_t \mid w_1, \dots, w_{t-1}\right),
\end{equation}
and is trained by minimising the token-level negative log-likelihood
$\mathcal{L}(\theta) = -\sum_{t=1}^{T} \log p_\theta(w_t \mid w_{<t})$. New
records are generated by sampling from this distribution and parsing the
emitted text back into typed feature values.

Tabula extends GReaT with modifications that target the efficiency
limitations of applying full natural-language models to tabular
data~\cite{Zhao2023Tabula}. It compresses the token representation of feature
names and categorical values to shorten sequences, applies a fixed padding
strategy, and considers initialising the network from random weights rather
than from a corpus-pretrained checkpoint, which removes the need for the model
to unlearn natural-language priors that are irrelevant to tabular
structure~\cite{Zhao2023Tabula}. Building on this framework,
\textcite{Miletic2024} serialise rows with token-sequence compression and
left padding, condition generation on the binary target variable as the
leading token, and benchmark the Pythia models against CTGAN, a conditional
tabular GAN that employs mode-specific normalisation~\cite{Xu2019CTGAN}. Using
a train-on-synthetic, test-on-real protocol with a random-forest estimator,
they report that the language models match or surpass the GAN baseline, with
utility increasing in model size, while the comparison of randomly initialised
against pretrained starting points was inconclusive. That study evaluates
utility alone and reports no privacy assessment, a gap that motivates the
present work.

\subsection{The Pythia Model Family}
\label{sec:bg:llm:pythia}

Pythia is a suite of decoder-only autoregressive transformers released to
support controlled study of how language models behave across scale and
training~\cite{Biderman2023Pythia}. The models span parameter counts from
tens of millions to twelve billion, and share a common architecture and a
single training corpus presented in an identical data order, so that
differences in behaviour can be attributed to scale rather than to confounded
data or optimisation choices~\cite{Biderman2023Pythia}. All checkpoints and
intermediate training states are publicly released, which makes the suite a
reproducible basis for fine-tuning experiments.

This thesis uses two members of the suite as scale anchors. Pythia-70M serves
as a compute-tractable baseline whose training and sampling costs are modest,
whereas Pythia-1B is the largest configuration examined by
\textcite{Miletic2024} and the scale at which the utility advantage of
language-model synthesis over tabular GANs is most pronounced. Restricting the
comparison to these two sizes isolates the effect of model capacity on the
privacy-utility trade-off while keeping the experimental budget feasible.

\subsection{LoRA Fine-tuning for Parameter-Efficient Adaptation}
\label{sec:bg:llm:lora}

Fine-tuning every parameter of a billion-parameter model is costly in memory
and storage, and is especially burdensome under differentially private
optimisation, which requires per-example gradient computation. Low-Rank
Adaptation (LoRA) reduces this cost by freezing the pretrained weight matrices
and learning a low-rank update in their place~\cite{Hu2022LoRA}. For a
pretrained weight matrix $W_0 \in \mathbb{R}^{d \times k}$, LoRA parameterises
the adapted layer as
\begin{equation}
\label{eq:bg:llm:lora}
h = W_0 x + \Delta W x = W_0 x + \frac{\alpha}{r}\, B A\, x,
\end{equation}
where $B \in \mathbb{R}^{d \times r}$, $A \in \mathbb{R}^{r \times k}$, the
rank $r \ll \min(d, k)$, and $\alpha$ is a fixed scaling constant. The matrix
$A$ is initialised from a random Gaussian distribution and $B$ is initialised
to zero, so that $\Delta W = 0$ at the start of training and the adapted model
initially coincides with the pretrained one~\cite{Hu2022LoRA}. Only $A$ and
$B$ are updated, which reduces the number of trainable parameters per layer
from $dk$ to $r(d + k)$, typically by several orders of magnitude.

This reduction is consequential for the differentially private fine-tuning
discussed in \S\ref{sec:bg:dp:dpsgd}. Because the Gaussian noise that DP-SGD
adds to enforce privacy is injected into every trainable coordinate,
restricting training to the low-dimensional LoRA parameters concentrates a
fixed privacy budget over far fewer dimensions and lowers the memory overhead
of per-example gradient clipping. Parameter-efficient adapters have been shown
to attain privacy-utility trade-offs competitive with, and often superior to,
full-model private fine-tuning~\cite{Yu2022DPFT,Li2022LargeDP}.

\subsection{DP-Augmented LLM Synthesis}
\label{sec:bg:llm:dp}

Language models are known to memorise and, under suitable prompting, to
reproduce verbatim training records, which makes uncontrolled LLM synthesis a
privacy hazard rather than a safeguard~\cite{Carlini2021Extracting}. Two
routes to constraining this leakage appear in the literature. The first
perturbs the generation process directly: SafeSynthDP pairs LLM-driven
synthesis with calibrated Laplace and Gaussian noise and traces the resulting
downstream accuracy against the privacy parameter $\varepsilon$, reporting
competitive performance and resistance to membership inference on a
text-classification benchmark~\cite{Nahid2024}. This output-side perturbation
does not, however, furnish a formal end-to-end training guarantee. The second
route enforces privacy during fine-tuning through DP-SGD
(\S\ref{sec:bg:dp:dpsgd}): \textcite{Yu2022DPFT} combine parameter-efficient
adapters with private optimisation to approach non-private utility on language
tasks at single-digit $\varepsilon$, and \textcite{Li2022LargeDP} report that
larger pretrained models act as stronger private learners, retaining more of
their utility as the privacy constraint tightens. This thesis adopts the
second route, fine-tuning Pythia-1B with DP-SGD via Opacus to instantiate the
Pythia-1B-DP synthesiser, and positions it against the non-private Pythia
variants and the GAN-based generators within a single privacy-utility
evaluation.

\section{Differential Privacy}
\label{sec:bg:dp}

\subsection{$(\varepsilon,\delta)$-Differential Privacy and Composition}
\label{sec:bg:dp:def}

Differential privacy provides a formal mathematical framework for quantifying
the privacy guarantees of randomised algorithms, and forms the theoretical
foundation for the differentially-private generators evaluated in this thesis.
A randomised mechanism $\mathcal{M}$ is said to satisfy
$(\varepsilon, \delta)$-differential privacy if, for every pair of neighbouring
datasets $D, D'$ that differ in a single record, and for every measurable
output subset $S \subseteq \mathrm{Range}(\mathcal{M})$:

\begin{equation}
\label{eq:bg:dp:def}
\Pr[\mathcal{M}(D) \in S] \;\le\; e^{\varepsilon}\,\Pr[\mathcal{M}(D') \in S] + \delta.
\end{equation}

The parameter $\varepsilon$ denotes the privacy budget: smaller values impose
stricter limits on the multiplicative ratio between output distributions, with
$\varepsilon = 0$ corresponding to statistical indistinguishability. The
relaxation term $\delta$ permits a small probability of catastrophic failure
and is conventionally chosen to be substantially smaller than the inverse
dataset size, with $\delta = 10^{-5}$ representing a standard convention for
datasets of order $10^5$ records~\cite{Dwork2014Foundations}.

The privacy guarantee of a deterministic query $f$ is converted into a
randomised mechanism by adding noise calibrated to the function's
\emph{global sensitivity}, defined as the maximum change in $f$ over all
pairs of neighbouring datasets. The Laplace and Gaussian mechanisms add noise
proportional to this sensitivity divided by $\varepsilon$ to satisfy pure
$\varepsilon$-DP and approximate $(\varepsilon, \delta)$-DP
respectively~\cite{Dwork2006DP,Dwork2014Foundations}.

Iterative procedures such as the differentially-private optimisation
algorithms employed in this thesis perform many noisy queries against the
private dataset, requiring composition rules that quantify how the cumulative
privacy cost grows. Basic sequential composition states that $k$ mechanisms,
each $(\varepsilon, 0)$-differentially private, yield a
$(k\varepsilon, 0)$-private composite, but this linear bound is loose for
practical training regimes that invoke thousands of gradient computations.
The advanced composition theorem of~\textcite{Dwork2014Foundations} tightens
the bound to approximately $O(\sqrt{k \log(1/\delta)}\,\varepsilon)$ for $k$
adaptive queries; even this remains insufficient for the repeated Gaussian
noise injections of DP-SGD, motivating the tighter accountants introduced in
§\ref{sec:bg:dp:rdp}.

\subsection{DP-SGD and Opacus}
\label{sec:bg:dp:dpsgd}

The Differentially-Private Stochastic Gradient Descent algorithm
of~\textcite{Abadi2016DPSGD} extends standard mini-batch SGD with two
privacy-preserving operations applied at every training step. Given a model
with parameters $\theta$, loss $\mathcal{L}$, clipping threshold $C$, noise
multiplier $\sigma$, and a Poisson-sampled mini-batch $B$ in which each
training example is included independently with probability $p$, the
algorithm computes per-example gradients
$g_i = \nabla_\theta \mathcal{L}(\theta, x_i)$ and clips each to a bounded
L2 norm:

\begin{equation}
\label{eq:bg:dp:clip}
\bar{g}_i = \frac{g_i}{\max\!\big(1,\; \lVert g_i \rVert_2 / C\big)}.
\end{equation}

The clipped per-example gradients are then summed, perturbed by isotropic
Gaussian noise, and averaged to obtain the noisy gradient used for the
parameter update:

\begin{equation}
\label{eq:bg:dp:update}
\tilde{g} \;=\; \frac{1}{\lvert B \rvert}\!\left(\sum_{i \in B} \bar{g}_i
   \;+\; \mathcal{N}\!\big(\mathbf{0},\, \sigma^{2} C^{2}\,\mathbf{I}\big)\right),
\qquad
\theta \;\leftarrow\; \theta - \eta\,\tilde{g}.
\end{equation}

Per-example clipping bounds the contribution any single training record can
make to a gradient update, capping the L2 sensitivity of the gradient sum at
$C$ and rendering the addition of Gaussian noise with standard deviation
$\sigma C$ sufficient to provide a formal privacy guarantee. The noise
multiplier $\sigma$ controls the privacy--utility trade-off: larger values
yield smaller privacy budgets at the cost of noisier gradients and slower
convergence. The total privacy expenditure across $T$ training steps is
computed using a tight accountant (§\ref{sec:bg:dp:rdp}) rather than by
naïvely composing the $(\varepsilon, \delta)$ guarantee of each step.

The Opacus library~\cite{Yousefpour2021Opacus} provides an implementation of
DP-SGD for the PyTorch ecosystem, exposing per-example gradient computation,
gradient clipping, noise injection, and privacy accounting as extensions to
standard training pipelines. Opacus is adopted in this thesis to train the
Pythia-1B-DP generator, with method-specific configuration details such as the
chosen noise multiplier and target $(\varepsilon, \delta)$ deferred to
Section~\ref{sec:method:generators}.

\subsection{The PATE Framework}
\label{sec:bg:dp:pate}

The Private Aggregation of Teacher Ensembles (PATE) framework
of~\textcite{Papernot2017PATE} offers an alternative route to
differentially-private learning that avoids the explicit gradient perturbation
of DP-SGD. The private training data is first partitioned into $N$ disjoint
subsets, and an independent \emph{teacher} model is trained on each partition
without any privacy mechanism. A separate \emph{student} model is then
trained on a corpus of unlabelled public data, with labels generated by
aggregating the teachers' predictions through a noisy voting mechanism:

\begin{equation}
\label{eq:bg:dp:pate}
\hat{y}(x) \;=\; \arg\max_{j}\;\Big\{\, n_{j}(x) \;+\; \mathrm{Lap}(1/\gamma)\,\Big\},
\end{equation}

where $n_j(x)$ denotes the number of teachers voting for class $j$ on input
$x$ and $\mathrm{Lap}(1/\gamma)$ is Laplace noise with scale $1/\gamma$
controlling the privacy cost per query.

The disjointness of the partitions is essential to the construction: because
any individual training record influences exactly one teacher, the
sensitivity of the vote count to the substitution of a single record is
bounded by one, allowing calibrated noise to confer a formal
$(\varepsilon, \delta)$-DP guarantee on the released labels. The student
model is trained exclusively on these noisy aggregated labels and never
accesses the private dataset, so the guarantee transfers to the trained
student via the post-processing property of differential privacy. A
favourable feature of the construction is that consensus among teachers
consumes a smaller portion of the privacy budget than disagreement, allowing
the framework to exploit easy queries efficiently. The Gaussian-noise
variant of PATE admits tighter analysis under Rényi differential privacy
(§\ref{sec:bg:dp:rdp}), and the application of the PATE construction to
generative adversarial networks is described in §\ref{sec:bg:gan:pategan}.

\subsection{Rényi DP and the Moments Accountant}
\label{sec:bg:dp:rdp}

The repeated application of the Gaussian mechanism in DP-SGD and the
iterative noisy aggregation in PATE both require composition tools more
refined than the basic and advanced composition theorems introduced in
§\ref{sec:bg:dp:def}. \textcite{Mironov2017RDP} proposed Rényi Differential
Privacy (RDP), a relaxation of standard differential privacy expressed in
terms of the Rényi divergence of order $\alpha > 1$. A randomised mechanism
$\mathcal{M}$ satisfies $(\alpha, \bar{\varepsilon})$-RDP if, for every pair
of neighbouring datasets $D, D'$:

\begin{equation}
\label{eq:bg:dp:rdp}
D_{\alpha}\!\big(\mathcal{M}(D)\,\Vert\,\mathcal{M}(D')\big)
\;\le\; \bar{\varepsilon},
\end{equation}

where $D_\alpha(P \Vert Q)$ denotes the Rényi divergence of order $\alpha$
between distributions $P$ and $Q$.

Rényi differential privacy admits a particularly simple composition rule:
the sequential application of $k$ mechanisms, each $(\alpha,
\bar{\varepsilon})$-RDP, yields a $(\alpha, k\bar{\varepsilon})$-RDP
composite. This linear growth at fixed $\alpha$ avoids the
square-root inflation incurred by the advanced composition theorem of
standard $(\varepsilon, \delta)$-DP. An $(\alpha, \bar{\varepsilon})$-RDP
mechanism may be converted to a standard $(\varepsilon, \delta)$-DP guarantee
via

\begin{equation}
\label{eq:bg:dp:conv}
\varepsilon \;=\; \bar{\varepsilon} \;+\; \frac{\log(1/\delta)}{\alpha - 1},
\end{equation}

for any $\delta \in (0, 1)$~\cite{Mironov2017RDP}. In practice the accountant
tracks the RDP cost at multiple orders $\alpha$ simultaneously during
training and, upon conversion, selects the order that minimises the resulting
$\varepsilon$ at the chosen $\delta$.

Rényi DP generalises the earlier \emph{moments accountant}
of~\textcite{Abadi2016DPSGD}, which tracked the higher moments of the
privacy-loss random variable to obtain tighter bounds for repeated Gaussian
queries than basic composition would yield. The Gaussian mechanism, the
subsampled Gaussian mechanism employed by DP-SGD with Poisson batching, and
the Laplace mechanism employed by PATE all admit closed-form RDP bounds at
every order $\alpha$, making RDP the standard accounting framework for both
algorithms.

\section{Evaluation Frameworks for Synthetic Data}
\label{sec:bg:eval}

% This thesis adopts a two-pillar evaluation framework: utility and privacy.
% The utility pillar is broad and itself contains two complementary aspects:
% statistical similarity (how well the synthetic distributions resemble the
% real ones) and predictive performance (how well downstream models trained
% on synthetic data behave). Privacy is assessed separately. The literature
% reviewed below populates both pillars.

\subsection{Utility -- Statistical Similarity to the Real Data}
\label{sec:bg:eval:similarity}

\textcite{Murtaza2023} propose a taxonomy that distinguishes univariate
metrics such as Kolmogorov--Smirnov (KS), Kullback--Leibler (KL) divergence,
and moment-based comparisons from multivariate metrics including correlation
matrix similarity, Maximum Mean Discrepancy, Jensen--Shannon divergence, and
Wasserstein distance, all used to quantify how closely synthetic distributions
match real data. \textcite{goncalves2020generation} complement these
distributional measures by combining KL divergence with pairwise correlation
differences (PCD) to assess both marginal and joint-distribution fidelity in
synthetic patient data. In the framework adopted here these metrics are
treated as one branch of utility, since a synthetic dataset that fails to
preserve the statistical structure of the real data cannot reasonably be
expected to support downstream tasks.

\subsection{Utility -- Predictive Performance of Downstream Models}
\label{sec:bg:eval:predictive}

\textcite{hernandez2022synthetic} survey tabular synthetic data generators and
emphasise model-based utility evaluations, including Train-on-Synthetic
Test-on-Real (TSTR) and related train--test regimes (TRTR/TRTS/TSTS), as
central tools for assessing whether synthetic data supports realistic
downstream inference. Reported metrics typically include accuracy, F1, and
AUC. \textcite{wang2019continuous}, \textcite{yang2019grouped},
\textcite{park2018data}, and \textcite{beaulieu2019privacy} all employ such
regimes; \textcite{emam2021optimizing} introduces Dimension-Wise Prediction
(DWP) as a complementary multivariate utility metric. Together with the
statistical-similarity metrics in the previous subsection, these
predictive-performance metrics constitute the second branch of the utility
pillar used throughout the thesis.

\subsection{Privacy Metrics: Distance, Inference, and Linkability}
\label{sec:bg:eval:privacy}

% Single combined section pulling together the privacy-metric paragraph from
% proposal §4 (the long paragraph beginning "Across the synthetic data
% literature").
Across the synthetic data literature, different studies employ diverse
privacy metrics to quantify disclosure risks. The evaluation framework by
\textcite{syninhospital} measures membership inference and attribute inference
for each patient record, together with privacy gain, which captures the
reduction of re-identification probability when disclosing synthetic rather
than real data. HealthGAN-based work evaluates privacy through
Nearest-Neighbour Adversarial Accuracy (AA) and privacy loss, and also
simulates membership inference attacks~\cite{yale2020synthesizing}.
\textcite{hernandez2022synthetic} summarise additional metrics, including
identity/attribute disclosure, Distance to Closest Record (DCR), membership
attack success rate, max real-to-synthetic similarity, formulation of
differential privacy, privacy loss, and Maximum Mean Discrepancy (MMD).
\textcite{Murtaza2023} systematically categorise privacy metrics into
integrated DP mechanisms (DP, Rényi DP, moments accountant), embedded
objectives such as $\varepsilon$-identifiability, $\ell$-diversity, and
probabilistic $k$-anonymity, and post-hoc metrics including nearest-neighbour
similarity, NN attribute estimation, reproduction rate, outlier similarity,
and distribution-based measures such as perplexity-distribution similarity.
Several papers also simulate membership inference by matching real and
synthetic records through distance thresholds (e.g., Hamming
distance~\cite{goncalves2020generation}) and evaluate attribute disclosure by
testing whether unknown attributes can be inferred from nearest synthetic
neighbours.

\section{Research Gap and Positioning}
\label{sec:bg:gap}

% TODO: Make explicit:
%  1. Few studies compare GAN- and LLM-based generators head-to-head on the
%     same medical benchmark with the same evaluation protocol.
%  2. Even fewer compare DP and non-DP variants \emph{within} the same family
%     (HealthGAN vs. PATE-GAN; Pythia vs. Pythia-DP).
%  3. Yale et al. \cite{healthgan} establish NNAA but do not benchmark against
%     LLMs; Miletic et al. \cite{Miletic2024} establish Pythia for tabular
%     data but skip privacy.
%  4. \textcite{hernandez2022synthetic} state that no single method or metric
%     is universally best — motivating the unified pipeline contributed here.
\textcite{hernandez2022synthetic}, in their systematic review of tabular
health data for synthetic data generation, conclude that no single method or
metric is universally best, and survey resemblance, privacy, and utility
metrics used by other authors. The thesis presented here directly addresses
this gap by applying a unified evaluation pipeline to four representative
generators spanning the GAN/LLM and DP/non-DP axes.

\cleardoublepage
